\newpage
\textbf{\huge CHAPTER  4}\\
\section{SOLUTION APPROACH AND ARCHITECTURE} \setcounter{section}{4}
\subsection{Introduction to the Proposed Solution}\setcounter{subsection}{1}
% \pagenumbering{arabic} 
\begin{justify}
    The critical bottlenecks identified in the legacy Apache Storm architecture—namely, the heavy reliance on external state management (HBase), the lack of exactly-once processing guarantees, and high network I/O latency—necessitated a fundamental paradigm shift in how Flipkart's Search team handles data ingestion. Under the guidance of my manager, Mr. Abhishek Agarwal, and mentor, Mr. Vinay, the core objective of my internship was formulated: to architect and implement a migration from Apache Storm to Apache Flink.

    Apache Flink is a robust, open-source stream processing framework designed for distributed, high-performing, and always-available data streaming applications. Unlike Storm, which processes events tuple-by-tuple without native state awareness, Flink treats state as a first-class citizen. This chapter details the architectural redesign, the implementation of Flink's stateful components, and the strategic approach taken to seamlessly transition the near real-time (NRT) ingestion pipeline.
\end{justify}

\vspace{0.25cm}

\subsection{Architectural Redesign with Apache Flink}
\begin{justify}
    The proposed architecture replaces the legacy Storm bolts with a highly optimized Flink DataStream topology. While the upstream message brokers (Varadhi, Kafka) and downstream sinks (Solr, Redis) remain consistent to ensure system compatibility, the intermediate processing engine undergoes a complete overhaul.
    
    
    
    \begin{itemize}
        \item \textbf{Source Connectors:} The new pipeline utilizes Flink's highly optimized Kafka Consumers to ingest the raw signal streams (Price, Availability, Catalog). These streams are immediately partitioned using Flink's \textit{keyBy()} operator based on the unique \textit{ProductID} or \textit{ListingID}. This ensures that all updates pertaining to a single product are deterministically routed to the same TaskManager node.
        \item \textbf{Data Enrichment via RichMapFunctions:} The pipeline leverages Flink's \textit{RichMapFunction} to handle external API calls necessary for base view transformations. To prevent these calls from becoming a new bottleneck, Flink's Asynchronous I/O API is implemented, allowing the system to handle multiple concurrent enrichment requests without blocking the main processing thread.
    \end{itemize}
\end{justify}

\subsection{Resolving the State Management Bottleneck (The Core Innovation)}
\begin{justify}
    The most significant architectural upgrade in this migration is the localization of state. In the Storm setup, every event required an external network call to the HBase EntityStore to fetch the current product state, merge the delta, and write it back.
    
    

    \begin{itemize}
        \item \textbf{Embedded RocksDB State Backend:} In the Flink architecture, we utilize RocksDB as the embedded state backend. RocksDB stores the state locally on the disk of the Flink TaskManager. When a price update arrives for a specific product, Flink retrieves the existing product entity directly from the local RocksDB instance in microseconds, applies the price delta, and updates the local state.
        \item \textbf{Eliminating External I/O:} By moving the "source of truth" into Flink's localized state during active processing, we effectively eliminate millions of external network calls to the HBase cluster per minute. HBase is now primarily used as a persistent backup and serving layer rather than a synchronous dependency for every stream event, drastically reducing overall ingestion latency.
    \end{itemize}
\end{justify}

\subsection{Handling Out-of-Order Events and Time Semantics}
\begin{justify}
    In distributed systems, network delays often cause events to arrive out of order (e.g., an older price update arriving after a newer one). Apache Storm struggled to natively handle event-time processing.
    
    To resolve this, the proposed solution implements Flink's \textbf{Event Time semantics} and \textbf{Watermarks}. Every ingested signal carries a timestamp of when it was generated. Flink uses Watermarks as a metric of progress in event time, allowing the system to temporarily buffer out-of-order events in customized Window functions. If a stale update arrives (its timestamp is older than the currently maintained state in RocksDB), the Flink \textit{ProcessFunction} simply discards it, ensuring that the final search index always reflects the most mathematically current state of the product.
\end{justify}

\subsection{Fault Tolerance and Exactly-Once Semantics}
\begin{justify}
    Flipkart's Search ecosystem cannot afford data loss or duplicate processing, as it directly impacts customer trust and revenue. The migration to Flink introduces strict "Exactly-Once" processing semantics through its distributed snapshotting mechanism, based on the Chandy-Lamport algorithm.
    
    \begin{itemize}
        \item \textbf{Checkpointing:} Flink periodically injects "Checkpoint Barriers" into the data stream. When these barriers flow through the operators, each operator saves a snapshot of its current state (from RocksDB) to a highly durable distributed file system (like HDFS). 
        \item \textbf{Failure Recovery:} In the event of a node failure or a pipeline crash, Flink automatically spins up a new TaskManager, restores the entire pipeline's state from the latest successful checkpoint, and rewinds the Kafka consumer offsets to the exact moment of the checkpoint. This guarantees that no product update is ever lost or applied twice, a massive improvement over Storm's "at-least-once" guarantees.
    \end{itemize}
\end{justify}

\subsection{Migration Strategy: Dual Writes and Shadow Testing}
\begin{justify}
    Replacing the core ingestion engine of a live, massive-scale e-commerce platform requires a meticulous rollout strategy to prevent any disruption to the live Search experience.

    The implementation is executed via a "Shadow Testing" approach. The Flink pipeline is deployed in parallel with the existing Storm topology. Both systems consume from the exact same Kafka source topics. However, the Flink pipeline initially writes its output to "Shadow Sinks" (isolated Kafka topics and a test Solr collection) rather than the production environment. We continuously run data validation scripts to compare the output entities generated by Flink against those generated by Storm. Once data parity is confirmed to be 100\% accurate across millions of test records, the Flink outputs will systematically replace Storm through a phased cutover.
\end{justify}

\subsection{Kafka Integration: FlinkKafkaConsumer and Exactly-Once Producer}
\begin{justify}
    A significant portion of my hands-on contribution during this internship was the design and implementation of the Kafka integration layer — the critical bridge between the upstream message bus and the Flink processing topology. While the broad migration strategy was defined at the team level, the precise configuration of the Kafka source and sink connectors, along with ensuring their correct interaction with Flink's checkpointing mechanism to achieve end-to-end exactly-once semantics, was my primary implementation responsibility.

    \textbf{FlinkKafkaConsumer (Source Integration).} The ingestion pipeline uses \texttt{FlinkKafkaConsumer} to read business signal events from Kafka. A key design decision was the configuration of the \textit{offset commit strategy}: rather than committing offsets to Kafka on every ACK (which would break exactly-once semantics in the event of a task failure mid-checkpoint), offsets are committed atomically as part of Flink's checkpoint completion. Specifically:
    \begin{itemize}
        \item The consumer is configured with \texttt{setStartFromGroupOffsets()} to resume correctly after restarts.
        \item \texttt{enable.auto.commit} is disabled; instead, the \texttt{FlinkKafkaConsumer} registers itself as a \textit{CheckpointListener} and commits the Kafka offsets corresponding to each checkpoint only after that checkpoint is confirmed complete by the JobManager.
        \item This guarantees that upon failure recovery, Flink restores its operator state from the checkpoint and simultaneously rewinds Kafka consumption to the exact offset committed at that checkpoint — achieving true end-to-end exactly-once processing.
    \end{itemize}

    \textbf{FlinkKafkaProducer (Sink Integration).} On the output side, processed product entities are written back to Kafka sink topics using \texttt{FlinkKafkaProducer} configured in \textit{exactly-once} mode, which leverages Kafka's transactional producer API (\texttt{transactional.id} + \texttt{initTransactions()} / \texttt{commitTransaction()} / \texttt{abortTransaction()}). This two-phase commit coordination — where Flink pre-commits a Kafka transaction as part of its checkpoint and finalizes it only upon checkpoint completion — ensures that downstream consumers of the sink topics never read duplicate or partial records, even in the presence of TaskManager failures.

    \textbf{Consumer Group and Partition Assignment.} Each Flink subtask is assigned a subset of Kafka partitions proportional to the pipeline's parallelism. The \textit{keyBy(listingId)} operator downstream ensures that all events for a given listing are co-located within a single Flink subtask, preserving ordering and enabling correct RocksDB state access without cross-node communication.

    \textbf{Challenges and Resolutions.}
    \begin{itemize}
        \item \textit{Transaction timeout alignment:} Kafka's default transaction timeout (15 minutes) must be greater than the maximum expected time between Flink checkpoints. I configured \texttt{transaction.max.timeout.ms} on the broker and \texttt{FlinkKafkaProducer}'s \texttt{kafkaProducerPoolSize} to prevent transaction expiry during prolonged backpressure events.
        \item \textit{Zombie fencing:} In the event of a network partition, a recovered producer could conflict with a zombie producer still holding a transaction. Kafka's epoch-based fencing mechanism, triggered automatically when a new transactional producer with the same \texttt{transactional.id} initialises, resolves this correctly.
    \end{itemize}
\end{justify}

\subsection{Latency Pill Integration: Observability Middleware}
\begin{justify}
    The second major contribution of my internship was integrating Flipkart's internal observability middleware — informally referred to as the \textit{Latency Pill} — into the Flink processing topology. This middleware is a standardised instrumentation library used across Flipkart's backend services to track, annotate, and report end-to-end processing latency at fine-grained stages of a data pipeline.

    \textbf{Purpose and Design.} The Latency Pill provides a lightweight, in-process tracing context that travels with each event as it passes through the system. Each stage of the pipeline — source consumption, enrichment, state merge, sink write — stamps its entry and exit timestamps into the Latency Pill context attached to the event. This accumulated trace is periodically flushed to an internal time-series metrics store (backed by an internal time-series database), enabling the Search team to monitor:
    \begin{itemize}
        \item \textbf{Per-stage latency breakdowns:} Precise isolation of which pipeline stage contributes most to end-to-end delay (e.g., enrichment API call vs. RocksDB state access vs. Kafka produce latency).
        \item \textbf{p50 / p95 / p99 latency dashboards:} Aggregated percentile dashboards used by on-call engineers to detect regressions and set SLA alerts.
        \item \textbf{Signal-type latency stratification:} Separate tracking for high-priority signals (out-of-stock, price drops) versus bulk catalog updates, enabling validation that priority streams are not starved by bulk traffic.
    \end{itemize}

    \textbf{Integration into the Flink Topology.} Integrating the Latency Pill into a Flink operator required careful attention to Flink's serialisation and state model:
    \begin{itemize}
        \item The Latency Pill context object was wrapped in a \texttt{TypeSerializer} compatible with Flink's managed memory model, ensuring it could be serialised correctly both during event transit and — critically — when snapshotted as part of Flink's checkpoint state.
        \item A custom \texttt{RichMapFunction} wrapper was implemented to inject Latency Pill stage stamps at the entry and exit of each operator, without modifying the core business logic of the operators themselves (open/closed principle).
        \item For the Kafka source, the pill's \texttt{originTimestamp} was initialised from the Kafka record's embedded event timestamp (set by the upstream publisher), ensuring that the measured latency represents the true end-to-end delay from business event generation to search index update — not merely the processing time within Flink.
    \end{itemize}

    \textbf{Validation and Impact.} Post-integration, the Latency Pill dashboards confirmed that Flink's localized RocksDB state access contributed less than 2ms at the p99 percentile to the per-event processing time, compared to the 80--150ms range observed for Storm's synchronous HBase reads under peak load. This quantitative visibility was instrumental in validating the performance benefits of the migration and in communicating results to senior engineering stakeholders.
\end{justify}